\section{The Divided Line}

\begin{quotex}
A man can be himself only so long as he is alone; and if he does not love solitude, he will not love freedom; for it is only when he is alone that he is really free.
\flright{\textsc{Arthur Schopenhauer}, \emph{Essays and Aphorisms}}
\end{quotex}


In the Republic, \textbf{Plato} describes four degrees of our apprehension of the world. From the highest to the lowest, there are:

\begin{enumerate}


\item \textbf{Intellection}. Grasping the Forms directly without the mediation of images or thoughts.

\item \textbf{Thought}. Our thoughts, ideas, discussions about the world.

\item \textbf{Belief}: Our opinions, including our knowledge of the world appearing to our senses.

\item \textbf{Imagery}. Created images of the world, fantasies.

\end{enumerate}


These are part of a continuum, not hard divisions, so they represent different degrees of understanding.

\paragraph{The Decision}


\begin{quotex}
If ... the soul departs from the body polluted and impure, in that it always is with the body and cares for and loves it and is bewitched by it, by desires and pleasures, so that it opines that nothing else is true except the corporeal, which one can grasp and see and drink and eat and use for sexual enjoyment, but is accustomed to hate and fear and flee from that which is dark and invisible to the eyes but is intelligible and grasped by philosophy, do you think that a soul like this will depart pure, itself by itself?
\flright{\textsc{Plato}}
\end{quotex}


What is described is more commonly expressed as ``Eat, drink, and be merry'' (Wisdom 8:15). This temptation arises from the seeming injustices in the world and it despairs in finding wisdom. While one should be grateful for having nourishment and moments of joy, to make that the only goal of life is to pollute the soul for death.

Yet what is dark and invisible to the senses, is intelligible to the philosopher. The philosopher is intimate with the Good, more intimate than simply ``seeing''. Plato describes it in terms of touching and nourishment, even erotic. The philosopher's soul ``touches'' divine reality and ``is nourished'' by the true and the divine; she even ``feasts on'' true being.

Whether to live in the shadows or to live in the light is the fundamental choice. Yet, choosing the latter may take you on a lonely path.

\paragraph{Measure, Number, and Weight}


\begin{quotex}
Thou hast ordered all things in measure and number and weight.
\flright{\textsc{Wisdom 11:20}}
\end{quotex}


\textbf{Roger Penrose}, who seems to be a natural Platonist, asserts in \textit{The Large, the Small, and the Human Mind}, that ``the entire physical world can, in principle, be described in terms of mathematics.'' This is consistent with the notion that Mathematics is the dividing line between the world of ideas and the world of senses. Also, therefore, with the notion that the rational is real and the real is rational, at least insofar as it applies to physics. Hence, the physical world can be understood by the Intellect provided the mathematics is understood.

The Book of Wisdom asserts the same thing.

\begin{itemize}


\item Measure refers to the shape of things, i.e., to Geometry;

\item Number to the count of things, i.e., to Arithmetic;

\item Weight to the mass of things.
\end{itemize}


We now know that Matter, or mass, is also convertible to Energy and to Frequency, so all is accounted for in that Triad: gravity, force, Maxwell's equations, and so on.

In the Phaedo, Plato speaks with approval of Anaxagoras' idea that ``the intellect is the orderer and cause of all things'', even in we might prefer to ascribe it to the \emph{Logos}. The cosmos must be good and beautiful, in that it is ordered by divine intelligence. That may be difficult sometimes to see, as the atheists are fond to point out.

\paragraph{Being and Privation}


To know something is to grasp the form behind the appearance. This is not dualism, since the form and the appearance are one. Actually, there cannot be anything without a form, or structure. But if forms are perfect, eternal, and unchanging, it would seem that things should also be perfect, eternal, and unchanging. Yet in the world, things do change and they are far from perfect.

I believe that Aristotle was trying to address this issue. If the forms are timeless and perfect, then why are not the instances of the forms perfect and unchanging? Aristotle points out that since things change, those things cannot fully embody the form. Rather there is something ``left over'', a surd, a lack, which he calls ``privation''. Hence things are a mix of the form (being) and privation (non-being). The rest follows from that initial observation.

However, in practice, it can be difficult to determine a thing's form. That is because the form is confused with its privation, which is the anti-form. That is why it has no being, although it is part of the world of our experience.

So if the form does not appear, then privation will be its substitute in the mind, even though it is the shadow, the doppelganger, or egregore of the thing. For physical objects, there may be less of a problem, but when the abstract things like political groupings, religions, families, etc., and even individual persona are mistaken for their shadows, confusion and even enmity arises. These are often accompanied with strong passions.

\paragraph{The Good}


Reality, as intelligible, consists of the forms that are revealed in things. The ``Good'', then, cannot be one of those forms that exists alongside other forms, hence it is beyond reality or Being. Things that are unified are Good. That is why things with privation will lack goodness. In particular, insofar as the human being is not unified within, he is not fully good. Our everyday experience should convince us that people are inconsistent from day to day, even hour to hour. They don't have a stable I or Self that unifies the disparate parts of the soul. The task then is to be made Whole, or Holy.

As for the totality of the world, perhaps God alone can grasp Reality in its completeness as a system of necessary relationships that plays out over time. Accept that on faith alone or strive to understand it.


\flrightit{Posted on 2020-05-03 by Cologero}

